\BourbakiExerciseGroup

\begin{enumerate}
\def\labelenumi{\arabic{enumi})}
\tightlist
\item
  \begin{enumerate}
  \def\labelenumii{\alph{enumii})}
  \tightlist
  \item
    Let \(q\) be a power of a prime number. Show that for every integer \(n > 1\) , if \(h_i\) ( \(1 \leqslant i \leqslant r\) ) are the distinct prime divisors of the integer \(n\) , then the number of elements \(\zeta \in \mathbf{F}_{q^n}\) such that \(\mathbf{F}_{q^n} = \mathbf{F}_q(\zeta)\) is equal to
  \end{enumerate}
\end{enumerate}

\[
\nu = q ^ {n} - \sum_ {i} q ^ {n / h _ {i}} + \sum_ {i <   j} q ^ {n / h _ {i} h _ {j}} - \sum_ {i <   j <   k} q ^ {n / h _ {i} h _ {j} h _ {k}} + \dots + (- 1) ^ {r} q ^ {n / h _ {1} h _ {2} \dots h _ {r}}
\]

(observe that such an element is characterized by the property of not belonging to any of the fields \(F_{q}^{n/h_{i}}\) ). If \(h_{1} \leqslant \cdots \leqslant h_{r}\) , we have

\[
q ^ {n} - \sum_ {i} q ^ {n / h _ {i}} \leqslant \nu \leqslant q ^ {n} - q ^ {n / h _ {1}}.
\]

Examine the case where \(n\) is a prime power.

\begin{enumerate}
\def\labelenumi{\alph{enumi})}
\setcounter{enumi}{1}
\item
  Let \(b_{n}(q)\) be the number of monic irreducible polynomials of degree \(n\) over \(\mathbf{F}_q\) . Deduce from \(a\) ) an expression for \(b_{n}(q)\) . In particular, if \(q \geqslant 3\) , show that the number of polynomials (not necessarily monic) of degree \(n\) in \(\mathbf{F}_q[\mathbf{X}]\) which are irreducible is at least equal to \(q^n (q - 2) / n \geqslant q^{n + 1} / 3n\) .
\item
  Prove that \(\sum_{d\mid n}db_d(q) = q^n\) (establish the identity \(1 / (1 - q\mathrm{T}) = \prod_{n}(1 - \mathrm{T}^n)^{-b_n(q)}\) ).
\end{enumerate}

* Using the Möbius inversion formula (Lie, II, p.~176, Appendix), show that

\[
n b _ {n} (q) = \sum_ {d \mid n} \mu (n / d) q ^ {d}.
\]

Compare this result with that obtained in \(b\) . *

\begin{enumerate}
\def\labelenumi{\alph{enumi})}
\setcounter{enumi}{3}
\tightlist
\item
  Let \(a_{n}(q)\) be the number of monic polynomials without multiple factors in \(\mathbf{F}_q[\mathbf{X}]\) . Show that \(a_0(q) = 1, \quad a_1(q) = q, \quad a_n(q) = q^n - q^{n-1}\) for \(n > 1\) (establish the identity \(q^n = \sum_i a_{n-2i}(q) q^i\) ).
\end{enumerate}

\begin{enumerate}
\def\labelenumi{\arabic{enumi})}
\setcounter{enumi}{1}
\item
  Show that the number of monic polynomials of degree \(n\) of \(\mathbf{F}_q[\mathbf{X}]\) which have no root in \(\mathbf{F}_q\) is \(q^{n - q}(q - 1)^q\) for \(n \geqslant q\) (observe that every function \(\mathbf{F}_q \to \mathbf{F}_q\) is a polynomial function).
\item
  Let \(\mathbf{K}\) be a finite field of \(q\) elements; for every prime number \(l\) let \(N_{l}\) be the union of the extensions of \(\mathbf{K}\) (contained in an algebraic closure \(\Omega\) of \(\mathbf{K}\) ) whose degree is a power of \(l\) . Show that \(\mathbf{N}_{l}\) is an abelian extension of \(\mathbf{K}\) whose Galois group is isomorphic to the group \(\mathbf{Z}_{l}\) of \(l\) -adic integers. Deduce that \(\Omega\) is isomorphic to the tensor product \(\otimes \mathbf{N}_{l}\) , and hence,
\end{enumerate}

that \(\operatorname{Gal}(\Omega / \mathbf{K})\) is isomorphic to the product \(\prod_{l} \mathbf{Z}_l\) .

* 4) Decompose the polynomial \(\mathbf{X}^4 + 1 \in \mathbf{F}_p[\mathbf{X}]\) into irreducible factors. *

\begin{enumerate}
\def\labelenumi{\arabic{enumi})}
\setcounter{enumi}{4}
\tightlist
\item
  Let \(\mathbf{K}\) be a finite field of \(q\) elements where \(q\) is not a power of 2, and let \(a_1, a_2, b\) be three elements of \(\mathbf{K}\) such that \(a_1a_2 \neq 0\) . Show that the number \(\nu\) of solutions \((x_1, x_2) \in \mathbf{K}^2\) of the equation \(a_1x_1^2 + a_2x_2^2 = b\) is given by the following formulae:
\end{enumerate}

\(1^{\circ}\) if \(b = 0\) and \(-a_{1}a_{2}\) is not a square in \(\mathbf{K},\nu = 1\)

\(2^{\circ}\) if \(b \neq 0\) and \(-a_{1}a_{2}\) is not a square in \(\mathbf{K}\) , \(\nu = q + 1\) ;

\(3^{\circ}\) if \(b = 0\) and \(-a_{1}a_{2}\) is a square in \(\mathbf{K}\) , \(\nu = 2q - 1\) ;

\(4^{\circ}\) if \(b \neq 0\) and \(-a_{1}a_{2}\) is a square in K, \(\nu = q - 1\) .

(When \(-a_{1}a_{2}\) is not a square, adjoin to K a root of \(\mathbf{X}^{2} + a_{1}a_{2}\) and use the fact that the multiplicative group of a finite field is cyclic.)

\begin{enumerate}
\def\labelenumi{\arabic{enumi})}
\setcounter{enumi}{5}
\tightlist
\item
  Let \(\mathbf{K}\) be a finite field of \(q\) elements where \(q\) is not a power of 2.
\end{enumerate}

\begin{enumerate}
\def\labelenumi{\alph{enumi})}
\tightlist
\item
  Let \(a_{1}, a_{2}, \ldots, a_{2m}, b\) be elements of K such that \(a_{1}a_{2} \ldots a_{2m} \neq 0\) . Show that the number \(\nu\) of solutions \((x_{1}, \ldots, x_{2m}) \in \mathbf{K}^{2m}\) of the equation \(a_{1}x_{1}^{2} + \cdots + a_{2m}x_{2m}^{2} = b\) is given by the following formulae:
\end{enumerate}

\(1^{\circ}\) if \(b = 0\) and \((-1)^{m}a_{1}a_{2}\ldots a_{2m}\) is not a square in K, \(\nu = q^{2m - 1} - q^{m} + q^{m - 1}\) ;

\(2^{\circ}\) if \(b \neq 0\) and \((-1)^{m}a_{1}a_{2} \ldots a_{2m}\) is not a square in K, \(\nu = q^{2m-1} + q^{m-1}\) ;

\(3^{\circ}\) if \(b = 0\) and \((-1)^{m}a_{1}a_{2}\ldots a_{2m}\) is a square in K, \(\nu = q^{2m - 1} + q^{m} - q^{m - 1}\) ;

\(4^{\circ}\) if \(b \neq 0\) and \((-1)^{m}a_{1}a_{2} \ldots a_{2m}\) is a square in K, \(\nu = q^{2m - 1} - q^{m - 1}\) .

(Argue by induction on \(m\) , using Ex. 5.)

\begin{enumerate}
\def\labelenumi{\alph{enumi})}
\setcounter{enumi}{1}
\tightlist
\item
  Let \(a_1, a_2, \ldots, a_{2m+1}\) , \(b\) be elements of \(K\) such that \(a_1a_2 \ldots a_{2m+1} \neq 0\) . Show that the number \(\nu\) of solutions \((x_1, x_2, \ldots, x_{2m+1}) \in K^{2m+1}\) of the equation
\end{enumerate}

\[
a _ {1} x _ {1} ^ {2} + \dots + a _ {2 m + 1} x _ {2 m + 1} ^ {2} = b
\]

is given by the following formulae :

\(1^{\circ}\) if \(b = 0, \nu = q^{2m}\) ;

\(2^{\circ}\) if \(b \neq 0\) and \((-1)^{m}a_{1}a_{2} \ldots a_{2m+1}b\) is not a square in K, \(\nu = q^{2m} - q^{m}\) ;

\(3^{\circ}\) if \(b \neq 0\) and \((-1)^{m}a_{1}a_{2} \ldots a_{2m+1}b\) is a square in \(\mathbf{K}\) , \(\nu = q^{2m} + q^{m}\) . (Reduce to the case \(a\) ).)

\begin{enumerate}
\def\labelenumi{\arabic{enumi})}
\setcounter{enumi}{6}
\tightlist
\item
  Let K be a finite field, E an algebraic extension of degree n of K and \((\alpha_{1}, \ldots, \alpha_{n})\) a normal basis of E over K; suppose that the cyclic Galois group \(\operatorname{Gal}(\operatorname{E}/\operatorname{K})\) is generated by the cyclic permutation \(\sigma = (\alpha_{1}\alpha_{2} \ldots \alpha_{n})\) . Let \(\tau\) be an automorphism of the vector K-space E such that for each \(x \in E\) , \(\tau(x)\) is conjugate to x over E (cf.~V, p.~154, Ex. 8).
\end{enumerate}

\begin{enumerate}
\def\labelenumi{\alph{enumi})}
\item
  Show that if \(\mathbf{K}\) has at least 3 elements, \(\tau\) is a K-automorphism of the field E (reduce to the case where \(\tau(\alpha_1) = \alpha_1\) and consider \(\tau(\alpha_1 + c\alpha_j)\) for \(c \in \mathbf{K}\) ).
\item
  If \(K = F_{2}\) , show that for each element \(\rho \in \operatorname{Gal}(E/K)\) , the relations \(\rho(\alpha_{r}) = \alpha_{r'}\) and \(\rho(\alpha_{s}) = \alpha_{s'}\) imply \(r' - r \equiv s' - s \pmod{n}\) . Deduce that if n \textgreater{} 5, \(\tau\) is a K-automorphism of the field E. (Note that if \(\tau(\alpha_{1}) = \alpha_{1}\) and \(\tau(\alpha_{1} + \alpha_{h}) = \alpha_{1} + \alpha_{k}\) with \(k \neq h\) , we must have \(h + k = n + 2\) ; deduce from this that necessarily \(\tau(\alpha_{2}) = \alpha_{n}\) and \(\tau(\alpha_{4}) = \alpha_{4}\) or \(\tau(\alpha_{4}) = \alpha_{n-2}\) which leads to a contradiction.)
\item
  For \(\mathbf{K} = \mathbf{F}_2\) and \(3 \leqslant n \leqslant 5\) the following automorphisms of the vector K-space E are not field automorphisms of E:
\end{enumerate}

\[
\text { for } \quad n = 3, \tau : (\alpha_ {1}, \alpha_ {2}, \alpha_ {3}) \mapsto (\alpha_ {1}, \alpha_ {3}, \alpha_ {2})
\]

\[
\text { for } \quad n = 4, \tau : (\alpha_ {1}, \alpha_ {2}, \alpha_ {3}, \alpha_ {4}) \mapsto (\alpha_ {1}, \alpha_ {4}, \alpha_ {3}, \alpha_ {2})
\]

\[
\text { for } \quad n = 5, \tau : (\alpha_ {1}, \alpha_ {2}, \alpha_ {3}, \alpha_ {4}, \alpha_ {5}) \mapsto (\alpha_ {1}, \alpha_ {5}, \alpha_ {4}, \alpha_ {3}, \alpha_ {2})
\]

but are such that \(\tau(x)\) is conjugate to \(x\) for all \(x \in E\) .

\begin{enumerate}
\def\labelenumi{\arabic{enumi})}
\setcounter{enumi}{7}
\tightlist
\item
  Let \(\mathbf{K}\) be a finite field of \(q\) elements and
\end{enumerate}

\[
\mathrm{P} (\mathbf {X}) = a _ {0} + a _ {1} \mathbf {X} + \dots + a _ {q - 2} \mathbf {X} ^ {q - 2}
\]

a polynomial of K{[}X{]} of degree q - 2. Let r be the rank of the matrix

\[
A = \left( \begin{array}{c c c c} a _ {0} & a _ {1} & \dots & a _ {q - 2} \\ a _ {1} & a _ {2} & \dots & a _ {0} \\ \dots & \dots & \dots & \dots \\ a _ {q - 2} & a _ {0} & \dots & a _ {q - 3} \end{array} \right).
\]

Show that the number of distinct roots of P in K, other than 0, is q - 1 - r. (Consider the product VA, where V is the Vandermonde matrix of elements \(x_{1}, x_{2}, \ldots, x_{q-1}\) of \(K^{*}\) .)

\begin{enumerate}
\def\labelenumi{\arabic{enumi})}
\setcounter{enumi}{8}
\tightlist
\item
  Let \(\mathbf{K}\) be a finite field of \(q\) elements and let \(r = q^n\) for an integer \(n \geqslant 1\) . For every polynomial \(\mathrm{P}(\mathbf{X}) = \sum_{j=0}^{m} a_j \mathbf{X}^j\) of \(\mathbf{K}[\mathbf{X}]\) , put
\end{enumerate}

\[
\hat {\mathrm{P}} (\mathrm{X}) = \sum_ {j} a _ {j} \mathrm{X} ^ {(r ^ {j} - 1) / (r - 1)}
\]

and

\[
\tilde {\mathrm{P}} (\mathrm{X}) = \mathrm{X} \hat {\mathrm{P}} \left(\mathrm{X} ^ {r - 1}\right) = \sum_ {j} a _ {j} \mathrm{X} ^ {r ^ {j}}.
\]

\begin{enumerate}
\def\labelenumi{\alph{enumi})}
\item
  Show that if P, Q are two polynomials in K{[}X{]} and if R = PQ, then we have \(\tilde{\mathbf{R}}(\mathbf{X}) = \tilde{\mathbf{P}}(\tilde{\mathbf{Q}}(\mathbf{X}))\) . Deduce that for two polynomials F, G of K{[}X{]} to be such that F divides G it is necessary and sufficient that \(\hat{F}\) should divide \(\hat{G}\) . (To see that the condition is sufficient, consider the Euclidean division of G by F.)
\item
  Let F be an irreducible polynomial in K{[}X{]} and let G be any polynomial in K{[}X{]}. If there exists an integer \(d \geqslant 1\) such that \(\hat{\mathrm{F}}(\mathbf{X}^{d})\) and \(\hat{\mathrm{G}}(\mathbf{X}^{d})\) have a common root, show that F divides G. (If H is the GCD of \(\hat{\mathrm{F}}(\mathbf{X}^{d})\) and \(\hat{\mathrm{G}}(\mathbf{X}^{d})\) consider the set of polynomials \(P \in K[X]\) such that H divides \(\hat{\mathrm{P}}(\mathbf{X}^{d})\) .)
\item
  Every polynomial \(\mathbf{F} \neq 0\) of \(\mathbf{K}[\mathbf{X}]\) divides a polynomial of the form \(\mathbf{X}^{\mathbf{N}} - 1\) . Show that if \(r - 1 = de\) and if \(\alpha\) is a root of \(\hat{\mathbf{F}}(\mathbf{X}^{d})\) , then \(\alpha\) belongs to the extension \(\mathbf{K}\) of degree \(n\mathbf{N}\) and deduce that the degree of every irreducible factor of \(\hat{\mathbf{F}}(\mathbf{X}^d)\) in \(\mathbf{K}[\mathbf{X}]\) divides \(n\mathbf{N}\) (use \(a\) )). If further \(e\) and \(d\mathbf{N}\) are relatively prime and if \(m\) is the degree of \(\alpha\) over the extension of \(\mathbf{K}\) of degree \(n\) , show that \(m\) divides \(\mathbf{N}\) (use the fact that \((r^{\mathbf{N}} - 1)/e \equiv d\mathbf{N} (\text{mod } e)\) ).
\end{enumerate}

\(d\) ) With the notation as in \(c\) ), suppose that \(\mathbf{F}\) is irreducible and that \(\mathbf{N}\) is the least of the integers \(h\) such that \(\mathbf{F}\) divides \(\mathbf{X}^h - 1\) . If \(e\) and \(d\mathbf{N}\) are relatively prime and if every prime factor of \(n\) divides \(\mathbf{N}\) , show that every irreducible factor of \(\hat{\mathbf{F}}(\mathbf{X}^d)\) is of degree \(n\mathbf{N}\) . (Note that in virtue of \(b\) ) we then have \(m = \mathbf{N}\) .)

\begin{enumerate}
\def\labelenumi{\alph{enumi})}
\setcounter{enumi}{4}
\tightlist
\item
  Deduce in particular from \(d\) ) that if \(\mathrm{P}(\mathbf{X}) = \sum_{j} a_j \mathbf{X}^j\) is an irreducible polynomial in
\end{enumerate}

K{[}X{]} and if N is the least of the integers h such that \(P(X)\) divides \(X^{h}-1\) , then every irreducible factor of \(\sum a_{j}X^{q^{j}-1}\) is of degree N.

\begin{enumerate}
\def\labelenumi{\arabic{enumi})}
\setcounter{enumi}{9}
\tightlist
\item
  Let \(t, n, k\) be three integers \(\geqslant 1\) such that \(t \leqslant n\) and \(n\) divides \(2^k - 1\) . Let \(S\) be the smallest set of integers such that \((1, t) \subset S \subset (1, n)\) and the condition \(\langle i \in S, 2i \equiv j (\text{mod } n) \text{ and } 1 \leqslant j \leqslant n \rangle\) implies \(j \in S\) . Put \(m = \text{Card } S\) . Let \(\xi\) be an element of order \(n\) in \(F_2^k\) and put \(P(X) = \prod_{j \in S} (X - \xi^j)\) .
\end{enumerate}

\begin{enumerate}
\def\labelenumi{\alph{enumi})}
\item
  Prove that \(\mathbf{P}(\mathbf{X})\in \mathbf{F}_2[\mathbf{X}]\) .
\item
  Show that if \(\mathbf{R}(\mathbf{X})\in \mathbf{F}_2[\mathbf{X}]\) is of degree \(< n\) and is a non-zero multiple of \(\mathrm{P}(\mathbf{X})\) , then there are strictly more than \(t\) non-zero coefficients.
\item
  Deduce the construction of a mapping \(\varphi: \mathbf{F}_2^n \to \mathbf{F}_2^m\) such that if \(u, v \in \mathbf{F}_2^n\) are distinct and have at most \(t\) distinct coordinates then \(\varphi(u) \equiv \varphi(v)\) .
\item
  Calculate S, m and P when t = 6, n = 15, k = 4.
\end{enumerate}

\begin{enumerate}
\def\labelenumi{\arabic{enumi})}
\setcounter{enumi}{10}
\tightlist
\item
  Let \(f\) be a monic polynomial in \(\mathbf{Z}[\mathbf{X}]\) , irreducible in \(\mathbf{Q}[\mathbf{X}]\) and let \(p\) be a prime number such that the canonical image \(\vec{f}\) of \(f\) in \(\mathbf{F}_p[\mathbf{X}]\) is a polynomial without multiple roots. If \(\Gamma\) (resp. \(\vec{\Gamma}\) ) is the Galois group of the polynomial \(f\) (resp. \(\vec{f}\) ) over \(\mathbf{Q}\) (resp. \(\mathbf{F}_p\) ), define an isomorphism of a subgroup of \(\Gamma\) onto \(\vec{\Gamma}\) . Deduce that if \(\vec{f} = g_1g_2\ldots g_r\) where the \(g_j\) are irreducible in \(\mathbf{F}_p[\mathbf{X}]\) and if \(g_j\) is of degree \(n_j\) (so that \(f\) has degree \(n = n_1 + n_2 + \dots + n_r\) ), then the group \(\Gamma\) , qua subgroup of \(\mathfrak{S}_n\) contains a product \(\sigma_1\sigma_2\ldots\sigma_r\) of cycles with disjoint supports, \(\sigma_j\) being of order \(n_j\) for \(1\leqslant j\leqslant r\) (cf.~I, p.~62).
\end{enumerate}

Generalize to the case where f is not monic, but where its leading coefficient is not divisible by p.

\begin{enumerate}
\def\labelenumi{\arabic{enumi})}
\setcounter{enumi}{11}
\item
  Let \(f\) be a polynomial of \(\mathbf{Z}[\mathbf{X}]\) of degree \(n\) and let \(p_1, p_2, p_3\) be three distinct prime numbers not dividing the leading coefficient of \(f\) ; let \(f_1, f_2, f_3\) be the canonical images of \(f\) in \(\mathbf{F}_{p_1}[\mathbf{X}], \mathbf{F}_{p_2}[\mathbf{X}], \mathbf{F}_{p_3}[\mathbf{X}]\) respectively. Suppose that \(f_1\) is the product of a linear factor and an irreducible factor of degree \(n - 1\) , \(f_2\) is the product of an irreducible factor of the second degree and irreducible factors of odd degrees, and finally \(f_3\) is irreducible. Show that under these conditions \(f\) is irreducible in \(\mathbf{Q}[\mathbf{X}]\) and that the Galois group of the polynomial \(f\) over \(\mathbf{Q}\) is isomorphic to the symmetric group \(\mathfrak{S}_n\) (« Dedekind's criterion »; use Ex. 11 as well as I, p.~63, Prop. 9).
\item
  Let \(p_1, p_2, \ldots, p_m\) be distinct odd prime numbers and \(P = p_1p_2 \ldots p_m\) ( \(m \geqslant 3\) ).
\end{enumerate}

\begin{enumerate}
\def\labelenumi{\alph{enumi})}
\tightlist
\item
  Given an integer \(n \geqslant 1\) , let \(E\) be the set of polynomials in \(\mathbf{Z}[\mathbf{X}]\) of degree \(\leqslant n\) , whose coefficients lie in the interval \([0, P[\) of \(\mathbf{N}\); we thus have \(\text{Card}(E) = P^{n+1}\) . Let \(k\) be a number such that \(0 < k < 1/3n\) ; show that among the polynomials \(f \in E\) there are at least \(k^m\mathbf{P}^{n + 1}\) such that for every index \(j\) such that \(1 \leqslant j \leqslant m\) , the canonical image \(f_j\) of \(f\) in \(\mathbf{F}_{p_j}[\mathbf{X}]\) is of degree \(n\) and irreducible (use Ex. 1, b)). Suppose moreover that \(k < \frac{1}{18(n - 2)}\) and \(n > 2\) ; arguing in the same way and using Dedekind's criterion (Ex. 12)
\end{enumerate}

show that there exists a set of polynomials \(f \in \mathbf{E}\) , of cardinal at least equal to \((1 - 3(1 - k)^m) \mathbf{P}^{n+1}\) which are of degree \(n\) , irreducible and whose Galois group over \(\mathbf{Q}\) is isomorphic to \(\mathfrak{S}_n\) .

\begin{enumerate}
\def\labelenumi{\alph{enumi})}
\setcounter{enumi}{1}
\tightlist
\item
  For every integer N let \(L_{n,N}\) be the set of polynomials of Z{[}X{]} of degree \(\leqslant n\) , whose coefficients belong to the interval \((-N, N)\) ; we thus have \(\text{Card}(L_{n,N}) = (2N + 1)^{n+1}\) . The integer n being fixed, show that for every \(\varepsilon\) such that \(0 < \varepsilon < 1\) , there exists an integer \(N_0\) such that for \(N \geqslant N_0\) , there are among the polynomials \(f \in L_{n,N}\) at least \((1 - \varepsilon)(2N + 1)^{n+1}\) which are irreducible of degree n, and whose Galois group over Q is isomorphic to \(S_n\) (use a)).
\end{enumerate}

* 14) Let K be a finite field, commutative or not, let Z be its centre, q the number of elements of Z and n the rank {[}K:Z{]}.

\begin{enumerate}
\def\labelenumi{\alph{enumi})}
\item
  If E is a subfield of K containing Z, show that \([E:Z]\) divides n.
\item
  Let \(x \in K\) be an element not belonging to Z; show that the number of distinct conjugates \(yxy^{-1}\) of x in \(K^{*}\) is of the form \((q^{n}-1)/(q^{d}-1)\) , where d divides n and is distinct from n (consider in K the set of elements permutable with x, and use a)).
\item
  Deduce from \(b\) ) that \(q - 1\) is divisible by the integer \(\Phi_n(q)\) (partition the group \(\mathbf{K}^*\) into conjugacy classes and use the relation (6) of \(V\) , p.~82).
\item
  Show that if \(n > 1\) , then \(\Phi_n(q) > (q - 1)^{\varphi(n)}\) (decompose \(\Phi_n(X)\) in the field \(C\) of complex numbers). Deduce that we must have \(K = Z\) , in other words, that \(K\) is necessarily commutative (Wedderburn's theorem). *
\end{enumerate}

